% Source: OpenAI2026AreaLaw, 06-transport.tex, transport:states, lines 393--401.
% Manuscript revision: adc7f1241b42e322a6451854ab7e4b4c146bf78a.
\section{Embedding the transport matrices}
\label{sec:transport_state_intertwine}

For positive definite inputs $(A_i)_{i\in I}$ and a weighted binary tree $T$,
write $M_T(A)$ for the root and $\Phi_{T,j}^A$ for its normalized leaf map.
For a root vector $v$ and $u\in\mathbb R$, define
\begin{align}
  \sigma_{T,j,u}^A(v)
    &=(\Phi_{T,j}^A)^*
       \bigl(|M_T(A)^{-iu}v\rangle\langle M_T(A)^{-iu}v|\bigr).
\end{align}
This definition uses the transport-state formula of~\cite{OpenAI2026AreaLaw},
display \texttt{transport:states}. The following auxiliary identities apply
to arbitrary root vectors under a common intertwiner of the inputs.

\begin{lemma}[Imaginary powers and a rectangular intertwiner]
  \label{lem:imaginary_power_intertwine}
  \lean{Matrix.Transport.imagPow_intertwine}
  \leanok
  Let $M,N$ be positive definite matrices on finite-dimensional complex
  Hilbert spaces and let $J$ map the second space to the first.
  If $MJ=JN$, then every $u\in\mathbb R$ satisfies
  \begin{align}
    M^{-iu}J&=JN^{-iu}.
  \end{align}
\end{lemma}

\begin{proof}
  \leanok
  The finite-dimensional functional calculus gives $(\log M)J=J\log N$.
  Multiplication by $-iu$ preserves this identity, as does taking matrix
  exponentials. Thus $\exp(-iu\log M)J=J\exp(-iu\log N)$.
\end{proof}

\begin{theorem}[Transport matrices under a common intertwiner]
  \label{thm:transport_state_intertwine}
  \lean{Matrix.Transport.transportState_intertwine}
  \leanok
  Let $(A_i)_{i\in I}$ and $(B_i)_{i\in I}$ be positive definite inputs
  on two finite-dimensional complex Hilbert spaces.
  Let $J$ map the second space to the first and suppose that
  $A_iJ=JB_i$ for every $i\in I$.
  Then every weighted binary tree $T$, label $j$, root vector $v$ on the
  second space and real $u$ satisfy
  \begin{align}
    \sigma_{T,j,u}^A(Jv)&=J\sigma_{T,j,u}^B(v)J^*.
  \end{align}
  The same tree, labels and weights are used on both sides. Repeated labels
  and zero terminal weights are allowed.
\end{theorem}

\begin{proof}
  \leanok
  \uses{thm:mean_tree_intertwine, lem:imaginary_power_intertwine,
    thm:trace_adjoint_leaf_map_intertwine}
  Theorem~\ref{thm:mean_tree_intertwine} gives $M_T(A)J=JM_T(B)$.
  Lemma~\ref{lem:imaginary_power_intertwine} therefore gives
  $M_T(A)^{-iu}Jv=JM_T(B)^{-iu}v$.
  Put $w=M_T(B)^{-iu}v$. Since $|Jw\rangle\langle Jw|
  =J|w\rangle\langle w|J^*$, Theorem~\ref{thm:trace_adjoint_leaf_map_intertwine}
  yields the stated identity.
\end{proof}

These identities hold for arbitrary rectangular $J$.
If $J$ is an isometry, it also preserves the norm of every root vector;
this is a separate assertion.
